\begin{abcliste}
\abc Die Formelsammlung (Geophysikalische Daten der Schweiz) gibt für Olten (\ang{47;21.1;} N, \SI{401}{\meter} ü.\,M.)
\begin{align}
 g &= \gol \approx \result{\golP{0}{4}}
\end{align}

\abc Mit Masse \MMO\ und Radius \rMO\ des Mars gilt
\begin{align}
 g_\mathrm{M} &= \gmaF \\
   &= \frac{\ncG\cdot\MM}{(\rM)^2} \\
   &= \gma \approx \result{\gmaP{0}{2}}
\end{align}

\abc Auf Meereshöhe hängt die Fallbeschleunigung von der geografischen Breite $\varphi = \brO$ ab (internationale Normalschwereformel):
\begin{align}
 g_0 &= \gnuF \\
   &= \gnu
\end{align}
In der Höhe $h = \hevO$ nimmt sie mit dem Quadrat des Abstands vom Erdmittelpunkt ab:
\begin{align}
 g &= \gevF \\
   &= \gnu\cdot\left(\frac{\rEar}{\rEar+\hev}\right)^2 \\
   &= \gev \approx \result{\gevP{0}{4}}
\end{align}
\end{abcliste}
